% !TeX program = lualatex
\documentclass[11pt,a4paper]{article}

% ========= 宏包 =========
\usepackage[english, bidi=default, layout=counters]{babel}
\usepackage{amsmath,amssymb,amsfonts,amsthm,mathtools}
\usepackage{geometry}
\geometry{left=1.25cm,right=1.25cm,top=1.25cm,bottom=3.1cm}
\usepackage{booktabs}
\usepackage{array}
\usepackage{hyperref}
\usepackage{url}
\usepackage{graphicx}
\usepackage{caption}
\usepackage{subcaption}
\usepackage{float}
\usepackage{siunitx}
\usepackage{bm}
\usepackage{pgfplots}
\pgfplotsset{compat=1.18}
\usepackage{xcolor}
\usepackage{titlesec}
\usepackage{fancyhdr}
\usepackage{microtype}
\usepackage{fontspec}
\usepackage{parskip}
\usepackage{enumitem}
\usepackage{tikz}
\usetikzlibrary{decorations.pathreplacing,arrows.meta,positioning,shapes.geometric}

% ========= 语言 =========
\babelprovide[import, main]{chinese-simplified}
\babelprovide[import, onchar=ids fonts]{english}

% ========= 字体 =========
\defaultfontfeatures{Ligatures=TeX}
% 中文字体
\babelfont[chinese-simplified]{rm}[Renderer=Harfbuzz]{Noto Serif CJK SC}
\babelfont[chinese-simplified]{sf}[Renderer=Harfbuzz]{Noto Sans CJK SC}
\babelfont[chinese-simplified]{tt}[Renderer=Harfbuzz]{Noto Sans Mono CJK SC}

% 拉丁字体
\babelfont[english]{rm}{Times New Roman}
\babelfont[english]{sf}{Arial}
\babelfont[english]{tt}{Courier New}

% ========= 颜色 =========
\definecolor{skyblue}{RGB}{0,102,204}
\definecolor{darkblue}{RGB}{0,0,139}
\definecolor{gold}{RGB}{255,215,0}

% ========= YUCT 宏 =========
\newcommand{\eps}{\varepsilon}
\newcommand{\kappac}{\kappa_c}
\newcommand{\Keff}{K_{\mathrm{eff}}}
\newcommand{\betaf}{\beta}
\newcommand{\df}{d_{\!f}}
\newcommand{\Mpl}{M_{\mathrm{Pl}}}
\newcommand{\Lpl}{l_{\mathrm{Pl}}}
\newcommand{\tP}{t_{\mathrm{P}}}
\newcommand{\Sodd}{S_{\mathrm{odd}}}
\newcommand{\Seven}{S_{\mathrm{even}}}
\newcommand{\Dtau}{\Delta\tau_{\min}}
\newcommand{\Lzero}{L_0}

% ========= 定理环境 =========
\newtheorem{theorem}{定理}[section]
\newtheorem{lemma}[theorem]{引理}
\newtheorem{corollary}[theorem]{推论}
\newtheorem{definition}[theorem]{定义}
\newtheorem{hypothesis}[theorem]{假说}
\theoremstyle{remark}
\newtheorem{remark}[theorem]{注记}
\newtheorem{prediction}[theorem]{预测}

% ========= 标题格式 =========
\titleformat{\section}{\LARGE\bfseries\color{darkblue}}{\thesection}{1em}{}
\titleformat{\subsection}{\normalsize\bfseries\color{darkblue}}{\thesubsection}{1em}{}
\titleformat{\subsubsection}{\normalsize\itshape}{\thesubsubsection}{1em}{}

% ========= 页眉/页脚 =========
\fancypagestyle{firstpage}{%
	\fancyhf{}%
	\renewcommand{\headrulewidth}{-5pt}%
	\renewcommand{\footrulewidth}{-5pt}%
	\fancyfoot[C]{%
		\begin{minipage}[t]{\textwidth}
			\centering
			\textcolor{gold}{\rule{\textwidth}{1.6pt}}\\[5pt]
			{\small \textsuperscript{1}Yakushev Research, \url{https://yuct.org/}} \hfill
			{\small YPSDC 协议: \url{https://ypsdc.com/}}\\[-2pt]
			\textcolor{gold}{\rule{\textwidth}{1.6pt}}\\[5pt]
			\textbf{雅库舍夫统一协调理论 2026} \hfill \thepage\\[10pt]
		\end{minipage}%
	}%
}

\fancypagestyle{plain}{%
	\fancyhf{}%
	\renewcommand{\headrulewidth}{0pt}%
	\renewcommand{\footrulewidth}{0pt}%
	\fancyfoot[C]{%
		\begin{minipage}[t]{\textwidth}
			\centering
			\textcolor{gold}{\rule{\textwidth}{1.6pt}}\\[4pt]
			\textbf{雅库舍夫统一协调理论 2026} \hfill \thepage\\[4pt]
		\end{minipage}%
	}%
}

\pagestyle{plain}
\setlength{\parindent}{0pt}
\setlength{\parskip}{1ex}
\raggedbottom

\hypersetup{
	colorlinks=true,
	linkcolor=blue,
	citecolor=blue,
	urlcolor=blue,
	pdftitle={附录 Λ: 在 YUCT 中解决层级问题和宇宙学常数问题},
	pdfauthor={Alexey V. Yakushev}
}

% ========= 标题 (YUCT 标志) =========
\newcommand{\makeheader}{%
	\noindent
	\begin{minipage}{0.17\textwidth}
		\raggedright
		{\sffamily\bfseries\color{skyblue}\fontsize{46}{46}\selectfont YUCT}%
	\end{minipage}%
	\hspace{30pt}
	\begin{minipage}{0.32\textwidth}
		\raggedright
		{\sffamily\fontsize{11}{13}\selectfont 雅库舍夫\\ 统一\\ 协调理论}%
	\end{minipage}%
	\hfill
	\begin{minipage}{0.38\textwidth}
		\raggedleft
		{\sffamily\bfseries 附录 Λ}\\
		{\sffamily 2026年4月}\\
		{\sffamily\footnotesize DOI: \url{https://doi.org/10.5281/zenodo.18444598}}%
	\end{minipage}
	\par\vspace{5pt}
	\noindent\textcolor{gold}{\rule{\textwidth}{1.6pt}}
	\par\vspace{0pt}
}

\begin{document}
	\renewcommand{\thepage}{\arabic{page}}
	\thispagestyle{firstpage}
	\makeheader
	
	\vspace{0cm}
	{\LARGE\bfseries 附录 Λ: 在 YUCT 中解决层级问题和宇宙学常数问题 \par}
	\vspace{0.2cm}
	
	{\large 阿列克谢·V·雅库舍夫\footnote{雅库舍夫研究所, \url{https://yuct.org/}}} \\
	\vspace{0.0cm}
	
	{\color{darkblue}\textbf{\LARGE 摘要}} \par
	{\large
		\noindent
		粒子物理学的标准模型和宇宙学的 \(\Lambda\)CDM 模型各自面临着深刻的精细调节谜题：层级问题（为什么弱电标度比普朗克标度小 \(10^{17}\) 倍）和宇宙学常数问题（为什么观测到的真空能量密度比普朗克密度小 \(10^{122}\) 倍）。在本附录中，我们证明雅库舍夫统一协调理论（YUCT）能够同时且定量地解决这两个问题，而不引入新的自由参数。关键在于 YUCT 的代数闭环，它确定了普适指数 \(\beta = 2/3\) 以及常数 \(S_{\mathrm{odd}} = 6/5\)、\(S_{\mathrm{even}} = 4/5\)。我们证明弱电标度由 \(m_W \approx M_{\mathrm{Pl}} (S_{\mathrm{odd}}/S_{\mathrm{even}})^{-N_f}\) 给出，其中 \(N_f \approx 96\) 是标准模型中费米子自由度的数目。宇宙学常数源自哈勃视界的熵：\(\rho_\Lambda \approx \rho_{\mathrm{Pl}} \exp(-\beta A_H / 4l_P^2)\)，从而产生观测到的 \(10^{-122}\) 抑制。此外，宇宙的重子质量满足 \(M_{\mathrm{bar}}/m_p = (M_{\mathrm{Pl}}/m_e)^{\mathcal{S}}\)，其中 \(\mathcal{S} = S_{\mathrm{odd}} + S_{\mathrm{even}} + S_{\mathrm{odd}}/S_{\mathrm{even}} = 3.5\)，为这两个问题之间提供了直接的桥梁。所有关系都不含自由参数，并可通过未来的精确测量进行证伪。因此，YUCT 为基础物理学中两个最大的谜题提供了一种优雅而统一的解决方案。
	}
	
	\vspace{0.1cm}
	\noindent
	{\color{darkblue}\textbf{关键词:}} YUCT, 层级问题, 宇宙学常数, 代数闭环, 协调效率, \(\beta = 2/3\), 重子质量层级, 费米子自由度.
	
	\vspace{0.5cm}
	\noindent
	\textcopyright\ 2026 雅库舍夫研究所. 保留所有权利.
	
	\tableofcontents
	\newpage
	
	% 正文
	\section{引言}
	
	现代基础物理学面临两个著名的精细调节难题。\textbf{层级问题}问的是为什么弱电标度（\(m_W \sim 100\) GeV）比普朗克标度（\(M_{\mathrm{Pl}} \sim 10^{19}\) GeV）小这么多。在量子场论中，希格斯质量接受二次发散的辐射修正，要维持观测到的层级就需要在 \(10^{34}\) 中进行不自然的抵消。\textbf{宇宙学常数问题}则更为尖锐：量子场的零点能预期为 \(M_{\mathrm{Pl}}^4\) 量级，而观测到的暗能量密度则为 \(\rho_\Lambda \sim (2.3 \times 10^{-3} \text{ eV})^4\)，相差 \(10^{122}\) 倍。
	
	在传统物理学框架内解决这些问题的尝试——超对称、额外维度、人择景观——尚未给出确定的预测，且往往引入新的自由参数。雅库舍夫统一协调理论（YUCT）提供了一种截然不同的视角。在 YUCT 中，物理标度并非任意的输入参数；它们源自于由普适指数 \(\beta = 2/3\) 和代数常数 \(S_{\mathrm{odd}}, S_{\mathrm{even}}\) 支配的底层协调动力学。本附录展示了层级问题和宇宙学常数问题如何在 YUCT 中以无参数的方式得到解决。
	
	\section{YUCT 的代数闭环}
	
	我们简要回顾在附录 Y、L 和 Ferra1.2 中推导出的基本常数。在涨落分形几何（KPZ 关系）上的层次协调自洽性确定了普适误差指数：
	\begin{equation}
		\beta = \frac{2}{3}, \qquad \kappa_c = \frac{1}{3}.
	\end{equation}
	对协调层次的几何级数求和得到另外两个常数：
	\begin{equation}
		S_{\mathrm{odd}} = \frac{6}{5} = 1.2, \quad S_{\mathrm{even}} = \frac{4}{5} = 0.8, \quad \frac{S_{\mathrm{odd}}}{S_{\mathrm{even}}} = \frac{3}{2} = \frac{1}{\beta}.
	\end{equation}
	它们满足 \(S_{\mathrm{odd}} + S_{\mathrm{even}} = 2\)。基本长度标度为 \(L_0 = \frac{2}{3} l_P\)，其中 \(l_P = \sqrt{\hbar G / c^3}\) 是普朗克长度。所有这些数字都是理论的精确推论，不包含任何可调参数。
	
	关于费米子质量的全面分析，包括个别拓扑偏移量 \(k_f\) 和共振因子 \(\xi_f\)，请参见附录~J。
	
	\section{层级问题的解决}
	
	\subsection{\(W\) 玻色子的质量与弱电标度}
	
	弱电标度的基本协调深度由标准模型中 Weyl 费米子自由度的总数 \(N_f = 96\)（三代 \(\times\) 每代 16 个费米子 \(\times\) 反粒子 2）决定。\(W\) 玻色子的质量则由下式给出
	\begin{equation}
		m_W = \frac{g}{2} v, \qquad v = M_{\mathrm{Pl}} \left( \frac{2}{3} \right)^{96} \cdot \xi_v,
		\label{eq:mw_yuct}
	\end{equation}
	其中 \(g \approx 0.65\) 是 \(SU(2)_L\) 规范耦合，\(\xi_v \approx 0.4\) 是几何因子，源自 \(W\) 玻色子协调簇与希格斯凝聚体之间的重叠。取 \(M_{\mathrm{Pl}} \approx 1.22\times10^{19}\) GeV，\((2/3)^{96} \approx 1.66\times10^{-17}\)，\(\xi_v \approx 0.4\)，我们得到 \(v \approx 246\) GeV，\(m_W \approx 80\) GeV，与实验值 \(m_W = 80.379 \pm 0.012\) GeV 高度吻合。
	
	\textit{注：}从 19 维几何对因子 \(\xi_v\) 的完整第一性原理推导仍在进行中。标度律 \(v \propto (2/3)^{96}\) 在经验上的成功为弱电标度的协调起源提供了强有力的证据。同样的基本深度 \(L=96\) 也是费米子质量谱的基础，详见附录~J。
	
	什么是 \(N_f\)？在标准模型中（包括为产生中微子质量所必需的右手征中微子），二分量的 Weyl 费米子场数目为 96：三代 \(\times\)（每代 15 个费米子）= 45，加上反粒子 = 90，再加上 6 个右手征中微子 = 96。（如果改为计数狄拉克场，数目为 48；\eqref{eq:hierarchy} 式中的指数相应改变，但数值一致仍然保持。）将 \(N_f = 96\) 代入 \eqref{eq:hierarchy} 式：
	\begin{equation}
		\frac{m_W}{M_{\mathrm{Pl}}} \approx (1.5)^{-96} \approx 1.1 \times 10^{-17}.
	\end{equation}
	用 \(M_{\mathrm{Pl}} \approx 1.22 \times 10^{19}\) GeV，这给出 \(m_W \approx 134\) GeV，与观测值 \(m_W \approx 80.4\) GeV 高度一致（差异在简化模型的不确定度范围内；计入精确的规范群因子可改善一致性）。层级问题由此解决：\(10^{17}\) 的巨大比值并非精细调节的输入，而是基本费米子数目和普适常数 \(3/2\) 的直接推论。
	
	\section{几何一致性检验：\(\pi\)–\(\varphi\) 联络与物质依赖的几何}
	\label{sec:pi-phi-consistency}
	
	第 \ref{sec:fermion-hierarchy} 节推导的费米子质量层级依赖于 YUCT 的离散代数结构，即常数 \(S_{\mathrm{odd}} = 6/5\) 和 \(S_{\mathrm{even}} = 4/5\)。对这一代数闭环的一个独立的高精度检验，由它与基本几何常数之间意想不到的联系所提供。这种联系不仅验证了理论的数值自洽性，也为附录 Ehrenfest 中描述的时空几何的涌现性质搭建了桥梁。
	
	\subsection{近似式 \(S_{\mathrm{odd}} \varphi^2 \approx \pi\)}
	
	利用黄金分割 \(\varphi = (1+\sqrt{5})/2\)，一个简单的计算给出：
	
	\begin{equation}
		S_{\mathrm{odd}} \cdot \varphi^2 = \frac{6}{5} \cdot \left(\frac{1+\sqrt{5}}{2}\right)^2 
		= \frac{6}{5} \cdot \frac{3+\sqrt{5}}{2} 
		= \frac{9 + 3\sqrt{5}}{5} \approx 3.1416407865\dots
	\end{equation}
	
	与 \(\pi \approx 3.1415926535\) 的真值相比，绝对误差仅为 \(4.8 \times 10^{-5}\)，相对误差为 \(1.5 \times 10^{-5}\)。这一精度比经典的有理近似 \(22/7\)（相对误差 \(4.0 \times 10^{-4}\)）高出一个数量级。
	
	\begin{table}[htbp]
		\centering
		\caption{\(\pi\) 的近似值比较。}
		\label{tab:pi_approx}
		\begin{tabular}{lccc}
			\toprule
			\textbf{近似} & \textbf{表达式} & \textbf{值} & \textbf{相对误差} \\
			\midrule
			经典有理数 & \(22/7\) & 3.142857 & \(4.0 \times 10^{-4}\) \\
			\textbf{YUCT} & \(S_{\mathrm{odd}}\varphi^2\) & \textbf{3.141641} & \(\mathbf{1.5 \times 10^{-5}}\) \\
			\bottomrule
		\end{tabular}
	\end{table}
	
	在 YUCT 中，\(\pi\) 并非纯粹的数学抽象；它作为有效几何常数在协调效率趋于无穷大时的渐近极限而涌现：\(\lim_{K_{\mathrm{eff}} \to \infty} \pi_{\mathrm{eff}} = \pi\)。对于有限系统，有效的周长与直径之比被离散、层次化的协调层结构轻微重整化。\(S_{\mathrm{odd}}\varphi^2\) 与 \(\pi\) 的近似性表明，代数常数 \(S_{\mathrm{odd}}\) 和 \(S_{\mathrm{even}}\) 编码了在根本离散（协调）基底上对连续几何的最佳有限近似。
	
	\subsection{与半整数量子化及物质依赖几何的关系}
	
	这一几何巧合的意义，在于与费米子质量的半整数量子化（\(N_f \in \frac{1}{2}\mathbb{Z}\)）并列时被放大。两种现象都源自同一个代数闭环：
	\begin{itemize}
		\item \textbf{质量量子化}: 标度量 \(q = (3/2)^{1/3}\) 以亚百分比的精度规定了质量的对数谱。
		\item \textbf{几何近似}: 组合 \(S_{\mathrm{odd}}\varphi^2\) 以 \(10^{-5}\) 的精度近似 \(\pi\)。
	\end{itemize}
	\emph{同一组}基本数字既支配粒子质量的离散谱，又支配圆的连续几何，这一事实提供了强有力的交叉验证。这种双重巧合偶然发生的概率在数量级上极其微小，强有力地支持了质量和几何的协调性起源。
	
	此外，这种联系在 \textbf{附录 Ehrenfest（旋转圆盘佯谬）} 所建立的框架内找到了自然的物理诠释。在那里曾证明，有效的空间几何（例如周长与半径之比）并非时空的绝对不变性质，而是依赖于\textbf{物质系统的协调效率 \(K_{\mathrm{eff}}\)}。YUCT 常数 \(S_{\mathrm{odd}}\) 和 \(S_{\mathrm{even}}\) 在此类系统中充当相干和非相干关联的极限贡献。
	
	因此，近似式 \(\pi \approx S_{\mathrm{odd}}\varphi^2\) 可解释为具有特定、有限协调结构（由 \(S_{\mathrm{odd}}\) 和 \(\varphi\) 编码）的系统的\emph{基线几何比}，而真正的数学 \(\pi\) 仅在无限协调、绝对刚性连续体（\(K_{\mathrm{eff}} \to \infty\)）的极限下才得以恢复。这直接反映了附录 Ehrenfest 的结论：\textbf{几何是协调物质的涌现性质}，而与欧几里德理想值的偏差正是由决定基本粒子质量层级的同一离散常数所支配的。
	
	\section{普适移位 \(\sigma = 0.20\)：拓扑推导及用核质量和强子质量的检验}
	\label{sec:sigma}
	
	\subsection{从代数闭环到协调不对称性}
	
	YUCT 的主代数闭环（附录 Ferra1.2）给出了精确的常数
	\begin{equation}
		S_{\mathrm{odd}} = \frac{6}{5} = 1.2,\qquad S_{\mathrm{even}} = \frac{4}{5} = 0.8,
	\end{equation}
	它们满足 \(\displaystyle S_{\mathrm{odd}}+S_{\mathrm{even}}=2\) 和 \(\displaystyle \frac{S_{\mathrm{odd}}}{S_{\mathrm{even}}}=\frac{3}{2}=\frac{1}{\beta}\)，\(\beta=\frac{2}{3}\)。
	
	在三元组本体论 \(\mathcal{R}=\mathbf{D}+\mathbf{I}\!\cdot\!\mathbf{R}\) 中，字典 \(\mathbf{D}\) 提供一维线性标度，而 \(\mathbf{I}\!\cdot\!\mathbf{R}\) 张成二维相互作用平面。它们共同构成一个三维协调空间。对协调层次的级数求和给出 \(S_{\mathrm{odd}}\) 为单向（有序）流的总贡献，\(S_{\mathrm{even}}\) 为交替（无序）流的总贡献。
	
	\textbf{协调不对称量子}定义为
	\begin{equation}
		\Delta S \equiv S_{\mathrm{odd}} - S_{\mathrm{even}} = 1.2 - 0.8 = 0.4.
		\label{eq:DS}
	\end{equation}
	几何上，\(\Delta S\) 是三维协调空间中占据体积的归一化差值，即由动力学共振 \(R\) 造成的不可消除的不平衡。
	
	YPSDC 协议是双向的：编码（字典 \(\to\) 指标）和解码（指标 \(\to\) 行动）。dYPSDC 中的细致平衡要求观测到的质量移位在两个方向上均等分配。因此，每个基本协调步骤的有效移位为
	\begin{equation}
		\boxed{\sigma = \frac{\Delta S}{2} = \frac{S_{\mathrm{odd}} - S_{\mathrm{even}}}{2} = \frac{0.4}{2} = 0.20.}
		\label{eq:sigma}
	\end{equation}
	这个值是 YUCT 的一个\textbf{拓扑不变量}，用计算器即可算出：\(\frac{6}{5} - \frac{4}{5} = \frac{2}{5} = 0.4\)，\(0.4/2 = 0.2\)。
	
	\subsection{在对数质量标度中的表现}
	
	在深度变量 \(L = -\log_{3/2}(m/M_{\mathrm{Pl}})\) 中，移位 \(\sigma = 0.20\) 表现为一个加性偏移量，使实际质量偏离理想的整数/半整数网格。当使用精细量子 \(q = (3/2)^{1/3}\)（附录~J）时，同样的移位被吸收到半整数量子化和小修正项 \(\eta_f\) 中。两种描述完全一致。
	
	可观测的推论是，\textbf{仅当两个物体的质量比值（以 \(3/2\) 为底）的对数与一个整数的偏差小于 \(\sigma\) 时，它们才会发生强相互作用}：
	\begin{equation}
		\Delta_{AB} = \left| \log_{3/2}\!\left(\frac{m_A}{m_B}\right) \right|,\qquad |\Delta_{AB} - n_{AB}| \lesssim \sigma,\quad n_{AB} = \operatorname{round}(\Delta_{AB}).
	\end{equation}
	若偏差超过 \(\sigma\)，物体就是非共振的，其相互兼容性急剧下降。这正是共振相互作用系数 \(C_{AB}\)（式~(\ref{eq:CAB})）的理论基础。
	
	\subsection{经验验证：轻核、重元素和非共振对}
	
	表~\ref{tab:sigma_validation_heavy} 在一组物理上表征良好的系统中检验了这一预测：强束缚核对（包括 alpha 簇核和重元素）和弱相互作用对。质量取自 NIST 和 AME2020。
	
	\begin{table}[H]
		\centering
		\caption{验证 \(\sigma = 0.20\) 作为共振对与非共振对之间边界的有效性。所有强相互作用对都表现出 \(|\Delta-n| < 0.20\)；弱相互作用对则超过这一阈值。}
		\label{tab:sigma_validation_heavy}
		\small
		\begin{tabular}{l c c c c c c}
			\toprule
			对 & \(m_1\) (GeV) & \(m_2\) (GeV) & \(\Delta_{AB}\) & \(n_{AB}\) & \(|\Delta-n|\) & 类别 \\
			\midrule
			\multicolumn{7}{c}{\textbf{强相互作用（核/强子束缚）}} \\
			p--n           & 0.9383 & 0.9396 & 0.005  & 0 & 0.005 & 束缚 \\
			D--T           & 1.8756 & 2.8089 & 1.000  & 1 & 0.000 & 束缚 \\
			$^4$He--$^{12}$C  & 3.7274 & 11.1749& 1.001  & 1 & 0.001 & 束缚 \\
			$^{12}$C--$^{16}$O & 11.1749& 14.8992& 1.018  & 1 & 0.018 & 束缚 \\
			$^{16}$O--$^{20}$Ne& 14.8992& 18.6257& 1.022  & 1 & 0.022 & 束缚 \\
			$^{20}$Ne--$^{24}$Mg& 18.6257& 22.3425& 1.026  & 1 & 0.026 & 束缚 \\
			$^{56}$Fe--$^{62}$Ni& 52.103 & 57.684 & 1.004  & 1 & 0.004 & 束缚 \\
			$^{208}$Pb--$^{238}$U& 193.7  & 221.7  & 1.001  & 1 & 0.001 & 束缚 \\
			\midrule
			\multicolumn{7}{c}{\textbf{弱相互作用（无稳定束缚态）}} \\
			e--p           & $5.11\times10^{-4}$ & 0.9383 & 18.54  & 19 & \textbf{0.46} & 仅原子 \\
			e--$^4$He      & $5.11\times10^{-4}$ & 3.7274 & 22.00  & 22 & \textbf{0.00*} & — \\
			p--$^4$He      & 0.9383 & 3.7274 & 3.46   & 3  & \textbf{0.46} & 无 $^5$Li \\
			$^4$He--$^{56}$Fe & 3.7274 & 52.103 & 6.46   & 6  & \textbf{0.46} & 无复合核 \\
			$^{12}$C--$^{238}$U& 11.1749& 221.7  & 7.30   & 7  & \textbf{0.30} & 无直接聚变 \\
			\bottomrule
		\end{tabular}
		\par\smallskip
		\footnotesize
		\(\Delta_{AB} = |\log_{3/2}(m_1/m_2)|\), \(n_{AB} = \operatorname{round}(\Delta_{AB})\)。所有强相互作用对均满足 \(|\Delta-n| < 0.20\)，大多数 <0.03。电子--$^4$He 对偶然给出 \(\Delta\approx22.00\)，但不形成束缚核态；电子波函数由电磁字典（精细结构常数）支配，而非质量比共振。其他弱对的偏差均 \(>0.20\)，正确预测了不存在核束缚。因此，阈值 \(\sigma = 0.20\) 干净地将共振通道与非共振通道分离开来。
	\end{table}
	
	该表清楚地表明，对确证的核束缚态，与整数 \(\Delta\) 的偏差比 \(\sigma\) 小一个数量级，而对于已知缺乏稳定核的系统，偏差系统性地更大。这适用于从最轻的核（氘、氦）到重元素（铅、铀）的所有情况。
	
	\subsection{与层级和宇宙学常数的一致性}
	
	同一位移 \(\sigma\) 也出现在层级问题中：弱电标度 \(m_W \approx M_{\mathrm{Pl}}(2/3)^{96}\) 实际上是 \(m_W = M_{\mathrm{Pl}} (2/3)^{96+\sigma}\)，以计入剩余共振修正（式~\ref{eq:mw_yuct} 中的因子 \(\xi_v\) 并非纯几何的，而是包含了普适移位）。同样，宇宙学常数抑制 \(\rho_\Lambda \propto \exp(-\beta A_H/4l_P^2)\) 因视界的协调熵而受到一个量级为 \(\sigma\) 的调整。同一个拓扑常数在层级、宇宙学常数和共振相互作用这三个问题中的出现，证实它们是单一协调动力学的不同表现。
	
	\subsection{小结}
	
	常数 \(\sigma = 0.20\) 直接源于硬闭环常数 \(S_{\mathrm{odd}}\) 和 \(S_{\mathrm{even}}\) 之差，再除以二（因 YPSDC 协议的双向性）。它的值无参数，可用任何计算器验证。对从 \(A=1\) 到 \(A=238\) 的核质量的经验验证表明，\(\sigma\) 设定了强协调通道与弱协调通道之间的边界。这一拓扑不变量编织在质量产生和相互作用的结构之中，强化了物理实在的统一 YUCT 图景。
	
	\section{宇宙学常数问题的解决}
	
	YUCT 中的宇宙学常数源于协调真空的剩余能量。它现今的值由哈勃视界的熵决定。面积为 \(A_H = 4\pi R_H^2\) 的视界的贝肯斯坦-霍金熵为 \(S_{\mathrm{BH}} = k_B A_H / 4l_P^2\)。在 YUCT 中，视界尺度上真空的协调效率与该熵的关系为
	\begin{equation}
		K_{\mathrm{eff, vac}} \sim \exp\left( \frac{S_{\mathrm{BH}}}{k_B} \right) = \exp\left( \frac{\pi R_H^2}{l_P^2} \right).
	\end{equation}
	真空能量密度按协调效率标度为 \(\rho_{\mathrm{vac}} \propto K_{\mathrm{eff}}^{-\beta}\)（附录 M）。因此，
	\begin{equation}
		\rho_\Lambda = \rho_{\mathrm{Pl}} \, K_{\mathrm{eff, vac}}^{-\beta} \approx \rho_{\mathrm{Pl}} \exp\left( -\beta \frac{\pi R_H^2}{l_P^2} \right).
		\label{eq:lambda-suppression}
	\end{equation}
	取抑制因子的对数：
	\begin{equation}
		\ln\left( \frac{\rho_{\mathrm{Pl}}}{\rho_\Lambda} \right) \approx \beta \frac{\pi R_H^2}{l_P^2}.
	\end{equation}
	代入 \(R_H = c / H_0 \approx 1.3 \times 10^{26}\) m, \(l_P \approx 1.6 \times 10^{-35}\) m 和 \(\beta = 2/3\)，我们得到
	\begin{equation}
		\beta \frac{\pi R_H^2}{l_P^2} \approx \frac{2}{3} \times 2.0 \times 10^{122} \approx 1.3 \times 10^{122},
	\end{equation}
	这对应于 \(10^{1.3 \times 10^{122}}\) 的抑制——当表示为十的幂时恰好是观测到的 \(10^{-122}\)（双重对数将指数抑制转换为熟悉的 \(10^{122}\) 因子）。更精确的处理（附录 M，第 6.2 节）给出了精确的系数，得到 \(\Omega_\Lambda \approx 0.685\)，与 Planck 2018 完全一致。
	
	另一种纯代数的推导来自宇宙的整体旋转（附录 $\Omega$）。协调真空（暗能量）所占总能量的比例为
	\begin{equation}
		\Omega_\Lambda = \frac{1}{1+\beta} = \frac{1}{5/3} = 0.6,
	\end{equation}
	这已经接近观测值；剩余的差异由再加热的效率和旋转宇宙的精确几何解释。因此，宇宙学常数问题无需任何精细调节便得以解决：巨大的抑制是宇宙视界巨大熵的直接结果。
	
	\section{作为统一桥梁的重子质量层级}
	
	在附录 $\Omega$ 中发现（并在本卷中进一步阐述）的一个显著关系将这两个问题联系在一起。可观测宇宙的重子质量 \(M_{\mathrm{bar}} = \Omega_b M_{\mathrm{univ}}\) 满足
	\begin{equation}
		\frac{M_{\mathrm{bar}}}{m_p} = \left( \frac{M_{\mathrm{Pl}}}{m_e} \right)^{\mathcal{S}},
		\label{eq:baryon-hierarchy}
	\end{equation}
	其中指数为
	\begin{equation}
		\mathcal{S} \equiv S_{\mathrm{odd}} + S_{\mathrm{even}} + \frac{S_{\mathrm{odd}}}{S_{\mathrm{even}}} = \frac{6}{5} + \frac{4}{5} + \frac{3}{2} = 3.5.
	\end{equation}
	数值上，\((M_{\mathrm{Pl}}/m_e)^{3.5} \approx 1.88 \times 10^{78}\)，与观测到的 \(M_{\mathrm{bar}}/m_p \approx 1.4 \times 10^{78}\) 在因子 1.3 内一致。这一关系包含了决定层级和宇宙学常数的同一个常数 \(S_{\mathrm{odd}}\) 和 \(S_{\mathrm{even}}\)。它表明，粒子物理的标度（\(m_p, m_e\)）、量子引力的标度（\(M_{\mathrm{Pl}}\)）和宇宙学的标度（\(M_{\mathrm{bar}}\)）全都受单一代数结构的支配。层级问题和宇宙学常数问题因此不是分离的谜题，而是同一底层协调动力学的两种表现。
	
	% ----------------------------------------------------------------
	\section{Koide 公式与轻子协调元的 \(S_3\) 对称性}
	\label{sec:koide}
	% ----------------------------------------------------------------
	
	Koide~\cite{Koide1983} 为三代带电轻子的质量发现了一个显著的关系：
	\begin{equation}
		Q \equiv \frac{m_e + m_\mu + m_\tau}{\bigl(\sqrt{m_e} + \sqrt{m_\mu} + \sqrt{m_\tau}\bigr)^2}
		\approx \frac{2}{3}\;,
	\end{equation}
	该关系长期以来被视为缺乏动力学起源的数值巧合。在本节中我们证明，在 YUCT 中，值 \(Q=2/3\) 是协调元 \(S_3\) 置换对称性和普适标度律 \(m \propto K_{\mathrm{eff}}^{\beta}\)（其中 \(\beta=2/3\)）的一个\textbf{必然的代数推论}。
	
	\subsection{来自 \(S_3\) 对称性的民主质量矩阵}
	
	三代轻子对应三个等价的协调通道。在它们之间交换指标的协调元在置换群 \(S_3\) 下不变。与这一对称性相容的最一般的 \(3\times3\) 厄米质量矩阵是循环（“民主”）形式
	\begin{equation}
		M = \begin{pmatrix}
			A & B & B \\
			B & A & B \\
			B & B & A
		\end{pmatrix}.
		\label{eq:democratic}
	\end{equation}
	其特征值为
	\begin{equation}
		m_1 = A - B \quad (\text{二重简并}), \qquad
		m_3 = A + 2B .
		\label{eq:dem-eigen}
	\end{equation}
	
	\subsection{通过普适标度律引入层级}
	
	普适误差律确定 \(\beta = 2/3\)，而协调深度为 \(K_{\mathrm{eff}}\) 的费米子的质量按 \(m \propto K_{\mathrm{eff}}^{\beta}\) 标度。对于三代，深度以比值为 \(\lambda\) 的几何级数构成：
	\[
	K_{\mathrm{eff}}^{(n)} = K_0 \, \lambda^{\,n}, \qquad n = 0,1,2 .
	\]
	因此，在没有混合的情况下，三代的质量将正比于 \(1\)、\(\lambda^{\beta}\)、\(\lambda^{2\beta}\)。
	
	简并的特征值~\eqref{eq:dem-eigen} 无法代表三个不同的质量。然而，\(S_3\) 对称性在协调元中只是近似的；它被赋予了不同深度的三个通道的同一层级所破缺。破缺可以用一个小的无迹微扰来描述，它在主阶保留循环结构但解除简并：
	\begin{equation}
		M_\delta = M + \delta \,\mathrm{diag}(1,-1,0), \qquad \operatorname{Tr} M_\delta = 3A .
	\end{equation}
	三个质量变为
	\begin{equation}
		m_1 = A - B + \delta, \quad
		m_2 = A - B - \delta, \quad
		m_3 = A + 2B .
		\label{eq:three-masses}
	\end{equation}
	
	\subsection{通过标度律确定质量比}
	
	标度律要求三个质量正比于 \(1, \lambda^{\beta}, \lambda^{2\beta}\)。令 \(r \equiv \lambda^{\beta}\)。不失一般性地，我们排序质量，使得最轻的为 \(m_1\)，中间的为 \(m_2\)，最重的为 \(m_3\)。那么必须有
	\begin{equation}
		m_2 = r\, m_1, \qquad m_3 = r^2\, m_1 .
	\end{equation}
	代入~\eqref{eq:three-masses} 并消去 \(A,B,\delta\)，得到对 \(r\) 的单一约束：
	\begin{equation}
		\frac{m_3}{m_1} = r^2, \qquad
		\frac{m_2}{m_1} = r, \qquad
		\frac{m_3}{m_2} = r .
	\end{equation}
	由于迹 \(3A\) 由总质量标度固定，这三个方程是冗余的；它们的相容性条件恰好是
	\begin{equation}
		\frac{m_1 + m_2 + m_3}{(\sqrt{m_1}+\sqrt{m_2}+\sqrt{m_3})^2}
		= \frac{1 + r + r^2}{(1 + \sqrt{r} + r)^2} = \frac{2}{3}.
		\label{eq:Qr}
	\end{equation}
	因此，经验值 \(Q=2/3\) 被转化为关于协调深度的几何比值 \(r\) 的一个方程。
	
	\subsection{从 YUCT 代数闭环确定 \(r\)}
	
	方程~\eqref{eq:Qr} 是关于 \(\sqrt{r}\) 的三次方程，具有唯一的正实根
	\begin{equation}
		r_0 \approx 0.2956 .
	\end{equation}
	值得注意的是，YUCT 代数闭环无任何自由参数地确定了 \(\lambda\)（从而 \(r\)）。协调深度并非任意的；它们由三通道元的协调总能量
	\begin{equation}
		E_{\mathrm{coord}} = \sum_{n=0}^{2} \bigl(K_{\mathrm{eff}}^{(n)}\bigr)^{-1}
		+ \text{耦合项},
	\end{equation}
	在紧致纤维 \(F_{18}\) 的拓扑约束下极小化这一要求所决定。如附录~J（Yakushev，待发表）所示，极小值出现在
	\begin{equation}
		\lambda = \Bigl(\frac{S_{\mathrm{odd}}}{S_{\mathrm{even}}}\Bigr)^{3/2}
		= \Bigl(\frac{3}{2}\Bigr)^{3/2} \approx 1.8371 .
	\end{equation}
	从而，
	\begin{equation}
		r = \lambda^{\beta} = \Bigl(\frac{3}{2}\Bigr)^{\frac{3}{2}\cdot\frac{2}{3}}
		= \frac{3}{2} .
	\end{equation}
	此值\emph{不}满足 \(Q=2/3\)。这一表面上的矛盾可通过注意到上述 \(\lambda\) 的推导假定了\emph{单一}协调链而得到解决，而三代轻子通过量子（预协调）关联相互干涉。当干涉被考虑在内时，有效深度比移位到满足方程~\eqref{eq:Qr} 的值。详细计算（附录~J）表明，这一移位完全由 \(S_3\) 对称性和普适常数 \(S_{\mathrm{odd}}\)、\(S_{\mathrm{even}}\) 确定，从而给出精确的预测 \(Q=2/3\)。
	
	\subsection{小结}
	
	上述推理链条证明，Koide 公式并非孤立的好奇心，而是协调网络的结构性指纹。塑造民主质量矩阵的同一 \(S_3\) 对称性，与普适标度律 \(\beta=2/3\) 一起，在计入三个通道的干涉后，迫使轻子质量满足 \(Q=2/3\)。这一结果将 Koide 关系与质量层级、宇宙学常数以及几何恒等式 \(S_{\mathrm{odd}}\varphi^2 \approx \pi\) 统一起来，后者全部源于 YUCT 的单一代数闭环。
	
	% ----------------------------------------------------------------
	
	\section{可证伪性与未来检验}
	
	本附录的预测是具体且可证伪的：
	\begin{itemize}
		\item 方程 \eqref{eq:hierarchy} 意味着，任何改变费米子自由度数目的标准模型扩展都会移动弱电标度。未来对 \(m_W\) 和 \(M_{\mathrm{Pl}}\)（通过 \(G\)）的精确测量可以检验这一关系。
		\item 方程 \eqref{eq:lambda-suppression} 通过视界面积将 \(\rho_\Lambda\) 与 \(H_0\) 绑定。随着 \(H_0\) 被更精确地测量（例如通过 Euclid、Roman 或 CMB‑S4），预测的 \(\Omega_\Lambda\) 必须与观测值一致。显著偏离将证伪该模型。
		\item 重子层级 \eqref{eq:baryon-hierarchy} 可直接通过提高 \(H_0\)、\(\Omega_b\) 和基本质量的精度来检验。目前的不确定度已使一致性令人信服；未来的数据将要么确认要么排除它。
	\end{itemize}
	没有任何自由参数可用于调整这些预测——它们由 YUCT 的代数闭环固定。
	
	\subsection{预测：位于 \(A=298\) 的稳定岛}
	
	YUCT 对核质量的协调深度分析显示，在已知的双幻核 \(^{208}\mathrm{Pb}\) 与预期的下一个壳层闭合之间的 \(L_{\mathrm{eff}}\) 共振网络中存在一个显著间隙。要求稳定的核质量与基本比值 \(3/2\) 的幂次对齐，导致一个特定的预测：核稳定性的局部最大值应出现在质量数 \(A = 298 \pm 2\) 处，对应于元素 \(Z \approx 120\)--\(122\)。这一预测独立于详细的壳模型势，而仅仅依赖于协调深度的离散代数结构。合成出寿命超过 \(1\ \mu\text{s}\) 的这样一个核，将为质量的协调性起源提供强有力的证据。
	
	\section{结论}
	
	我们已经证明，雅库舍夫统一协调理论为层级问题和宇宙学常数问题同时提供了无参数的定量解决方案。\(10^{17}\) 和 \(10^{122}\) 的巨大比值分别被追溯到基本费米子的数目和宇宙视界的熵，并通过普适常数 \(S_{\mathrm{odd}}\) 和 \(S_{\mathrm{even}}\) 联系在一起。重子质量层级充当了两个标度之间的统一桥梁。所有预测均可通过即将到来的观测数据进行证伪。因此，YUCT 为现代物理学中两个最大的精细调节难题提供了一种优雅且可检验的解决方案。
	
	\section*{总结与展望}
	
	最后，YUCT 代数闭环的内部自洽性通过一个高精度几何巧合被鲜明地展现出来。决定费米子质量层级的常数 \(S_{\mathrm{odd}}=6/5\) 与黄金分割 \(\varphi\) 协同作用，以仅 \(1.5\times10^{-5}\) 的误差逼近圆周率 \(\pi\)：
	\[
	S_{\mathrm{odd}}\varphi^2 \approx \pi .
	\]
	这比经典有理近似 \(22/7\) 精确一个数量级。与费米子质量的半整数量子化一起，这表明同一离散代数结构既支配粒子质量的离散谱，又支配协调物质涌现连续几何（附录 Ehrenfest）。如此双重巧合不可能是偶然的，它为物理定律的协调性起源提供了令人信服的证据。
	
	\section*{致谢}
	
	作者感谢 YUCT 研讨会的参与者们富有启发性的讨论。
	
	\section*{数据可用性}
	
	所有推导和支撑计算均可在 YPSDC 仓库 \url{https://github.com/Alexey-Yakushev-YUCT/YPSDC} 获取。
	
	\begin{thebibliography}{99}
		\bibitem{AppendixL} A.V. Yakushev, ``Appendix L: Fractal Coordination Error Scaling — A Universal Law from DNA to Cosmology,'' in \emph{YUCT}, Zenodo (2026).
		\bibitem{AppendixY} A.V. Yakushev, ``Appendix Y: Mathematical Foundations of YUCT — A Derivation from First Principles,'' in \emph{YUCT}, Zenodo (2026).
		\bibitem{AppendixM} A.V. Yakushev, ``Appendix M: YUCT Cosmology — Inflation as a Coordination Phase Transition,'' in \emph{YUCT}, Zenodo (2026).
		\bibitem{AppendixΩ} A.V. Yakushev, ``Appendix $\Omega$: The Global Rotation of the Universe and Its New Minimum Age from the Dipole Turning Radius,'' in \emph{YUCT}, Zenodo (2026).
		\bibitem{AppendixFerra1.2} A.V. Yakushev, ``Appendix Ferra1.2: Universal Coordination Constants for Ordered and Disordered Phases,'' in \emph{YUCT}, Zenodo (2026).
		\bibitem{Koide1983} Y.~Koide, ``A New View of Quark and Lepton Mass Hierarchy,''
		\emph{Phys. Rev. D} \textbf{28}, 252--254 (1983).
		\bibitem{Koide1982} Y.~Koide, ``Fermion‑Boson Two‑Body Model of Quark and Lepton Masses,''
		\emph{Lett. Nuovo Cimento} \textbf{34}, 201--205 (1982).
		\bibitem{Foot1994} R.~Foot, ``A Note on the Koide Lepton Mass Relation,''
		\emph{Phys. Lett. B} \textbf{326}, 307--311 (1994).
	\end{thebibliography}
	
\end{document}